\section{Respuesta a la RQ3}
\label{sec:limites:respuesta}

La RQ3 pregunta cómo se relacionan los límites de uso y los \textit{rates} en las APIs de correo, y qué implicaciones tiene esa relación para el consumidor.

En seis de las ocho APIs (Postmark, Resend, SendGrid, MailerSend, Brevo y Mailgun) la cuota es la restricción activa, el límite que se alcanza primero y frena el envío (Sección~\ref{sec:limites:geometrica}).
En lo que el \textit{rate} admite durante la ventana cabe entre 103,7 veces la cuota (MailerSend Starter) y 129.600.000 veces (el plan Free de Postmark, cuyo \textit{rate} no está publicado y sale de la cabecera que devuelve la API).
La condición es que el \textit{rate} y la cuota tengan valor en la Datasheet, publicado o medido, y que la cuenta no esté en revisión, en pruebas ni restringida por reputación.
Con esa condición, quien contrate una de las seis puede dimensionar el volumen sin consultar el \textit{rate}.

El cociente $R \cdot W / Q$ baja a 2,59 en MailerSend Professional 2M, un volumen contratado que la Datasheet no recoge, y a 0,43 en Starter 1M con la cuenta en revisión, el único caso de las dos tablas con cociente calculable en que manda el \textit{rate}.
En Amazon SES y Mailchimp Transactional no se puede decir que mande la cuota.
Amazon SES no publica la cuota ni el caudal de producción, así que la relación no se puede calcular antes de sacar la cuenta del entorno de pruebas.
En Mailchimp Transactional la restricción activa es el caudal: la cuenta de prueba tiene una franquicia de 500 correos y un tope de 25 correos por hora, de modo que agotarla lleva 20 horas.

Las implicaciones para quien contrata salen de tres casos en que la cuota no basta para dimensionar el envío.
El primero es el estado de la cuenta, que condiciona la capacidad en seis de las ocho APIs sin que ninguna publique el umbral que dispara el cambio.
Solo MailerSend da una cifra para cada estado, 120 peticiones por minuto con la cuenta sana frente a 10 en revisión.

El segundo caso es la distancia entre el \textit{rate} publicado y el observado.
Solo coinciden en cuatro APIs, las que escriben a qué se aplica el valor.
Resend lo fija para cada equipo de la cuenta, MailerSend para cada estado de la cuenta, Brevo para cada \textit{endpoint} en su categoría General y Mailchimp Transactional para la cuenta de prueba.
En Mailgun, en cambio, no coinciden las 1.500 peticiones por segundo y por IP de sus preguntas frecuentes y el contador de la cabecera del envío, 80.000 peticiones por ventana de unos diez segundos.

Amazon SES y Mailchimp Transactional no emiten ninguna cabecera de caudal, y Mailgun solo la devuelve tras un envío aceptado.
Para comprobar antes de enviar cuánta capacidad queda hace falta un \textit{endpoint} que la informe, y solo la mitad del catálogo lo tiene.
Amazon SES, Mailchimp Transactional y MailerSend devuelven el límite, y Brevo solo el saldo de correos.

El tercer caso es el factor $k$, los destinatarios por petición, que pasa el \textit{rate} de peticiones a correos.
Solo vale si cada destinatario descuenta un correo de la cuota, y eso solo consta en Resend, SendGrid, Brevo y Amazon SES.
En Postmark, MailerSend, Mailgun y Mailchimp Transactional el paso de peticiones a correos es una suposición.

El contraste con PRIME4API de la Tabla~\ref{tab:a06-presupuesto}, que da el tiempo mínimo para enviar 100.000 correos con y sin presupuesto, añade dos implicaciones más.
La primera es que la Datasheet no dice cuál de los dos regímenes aplica el proveedor: en Postmark Basic, la respuesta es de 270 días si la cuota es un corte estricto y de 19 segundos si se paga el sobrecoste.
La segunda es que el motor lee como infinito el caudal que la Datasheet deja vacío: con ese bloque vacío, agotaba en 0 segundos los 10.000 correos de Postmark Basic, cuyo \textit{rate} el proveedor no publica.

En el borde de la cuota hay tres comportamientos, que recoge la Sección~\ref{sec:limites:catalogo} y no el contraste con el motor.
Amazon SES rechaza la llamada y descarta el mensaje, y SendGrid, MailerSend y Mailchimp Transactional detienen el envío donde no hay sobrecoste.
Brevo retiene en cola hasta 1.000 correos por encima de la cuota diaria y no entrega los que pasan de esa cifra, que es lo que modela la propuesta P-09 del Apartado~\ref{subsec:evolucion:p09}.

Esta respuesta descansa sobre una base estrecha.
El par formado por el \textit{rate} y la cuota solo tiene valor en la Datasheet de seis APIs.
En cuatro de ellas el proveedor publica el \textit{rate} del envío, y en Postmark y Mailgun sale de una cabecera medida: en Mailgun, 80.000 peticiones por una ventana de unos diez segundos, que la Tabla~\ref{tab:limites-rates-b} divide entre diez.
La cifra que Mailgun publica, 1.500 peticiones por segundo y por IP, tiene otro ámbito.
De los cuatro \textit{rates} publicados, el de SendGrid no se liga a ningún plan y el de Resend es el valor por defecto del equipo.
El de MailerSend solo vale para una cuenta aprobada, y la página de Brevo se contradice sobre qué planes incluyen cada una de sus tres categorías de caudal.

Los valores medidos tampoco describen el plan contratado: en MailerSend, Amazon SES y Mailchimp Transactional corresponden a cuentas sin aprobar o en pruebas, y en Postmark y Mailgun a una sola cuenta gratuita.
La relación entre cuota y \textit{rate} queda establecida en seis APIs, pero el margen exacto entre ambos depende de valores que los proveedores no fijan por escrito.
